% !TEX root = ../../ctfp-print.tex

\lettrine[lhang=0.17]{看}{起来在范畴论中}一切事物都彼此关联, 一切事物都能从许多角度来看. 举个例子,
\hyperref[products-and-coproducts]{积}的普遍构造. 既然我们已经了解了更多 \hyperref[functors]{函子}
和 \hyperref[natural-transformations]{自然转换} 的知识, 能不能把它简化, 乃至推广呢? 试试看.

\begin{figure}[H]
  \centering
  \includegraphics[width=0.3\textwidth]{images/productpattern.jpg}
\end{figure}

\noindent
积的构造始于挑选两个物件 $a$ 和 $b$, 我们要构造它们的积. 但\emph{挑选物件}是什么意思?
我们能不能用更范畴论的说法来重述这个动作? 两个物件构成一个模式 --- 一个非常简单的模式.
我们可以把这个模式抽象成一个范畴 --- 一个非常简单的范畴, 但它毕竟是个范畴. 我们叫它
$\cat{2}$. 它只有两个物件 $1$ 和 $2$, 除了两个必然的恒等态射之外没有别的态射. 现在我们可以把
``在 $\cat{C}$ 中挑选两个物件''重述成: 定义一个从范畴 $\cat{2}$ 到 $\cat{C}$ 的函子 $D$.
函子把物件映到物件, 所以它的像就是两个物件 (也可能是一个, 如果函子把物件合并掉了, 那也没关系).
它也映射态射 --- 在这里它只是把恒等态射映到恒等态射.

这个进路的好处在于, 它建立在范畴论的概念之上, 从而避开了诸如 ``挑选物件'' 这类含糊的说法 ---
那简直是我们狩猎采集祖先的词汇里直接搬来的. 顺带一提, 它也很容易推广, 因为没什么能阻止我们用比
$\cat{2}$ 更复杂的范畴来定义自己的模式.

\begin{figure}[H]
  \centering
  \includegraphics[width=0.35\textwidth]{images/two.jpg}
\end{figure}

\noindent
接着往下. 积的定义中的下一步是挑选候选物件 $c$. 这里同样可以借某个来自单例范畴的函子来重述这个
挑选动作. 的确, 如果我们会用 Kan 扩张, 那就该这么做. 但既然我们还没准备好谈 Kan 扩张, 有个
别的技巧可用: 从同一个范畴 $\cat{2}$ 到 $\cat{C}$ 的常数函子 $\Delta$. 在 $\cat{C}$ 中挑选
$c$ 可以靠 $\Delta_c$ 完成. 记住, $\Delta_c$ 把所有物件都映到 $c$, 把所有态射都映到
$\idarrow[c]$.

\begin{figure}[H]
  \centering
  \includegraphics[width=0.35\textwidth]{images/twodelta.jpg}
\end{figure}

\noindent
现在我们有两个函子, $\Delta_c$ 和 $D$, 都跑在 $\cat{2}$ 和 $\cat{C}$ 之间, 那么很自然地要问
它们之间的自然转换. 由于 $\cat{2}$ 里只有两个物件, 一个自然转换会有两个分量. $\cat{2}$ 中的
物件 $1$ 被 $\Delta_c$ 映到 $c$, 被 $D$ 映到 $a$. 所以 $\Delta_c$ 和 $D$ 之间的自然转换在
$1$ 处的分量是一个从 $c$ 到 $a$ 的态射, 我们可以叫它 $p$. 同理, 第二个分量是从 $c$ 到 $b$ 的
态射 $q$ --- 即 $\cat{2}$ 中的物件 $2$ 在 $D$ 下的像. 可这些恰恰就是我们最初定义积时用到的
两个投影. 所以与其谈论挑选物件和投影, 我们完全可以只谈挑选函子和自然转换. 恰好在这个简单情形里,
我们那个转换的自然性条件是平凡满足的, 因为 $\cat{2}$ 里没有 (恒等之外的) 态射.

\begin{figure}[H]
  \centering
  \includegraphics[width=0.35\textwidth]{images/productcone.jpg}
\end{figure}

\noindent
把这个构造推广到 $\cat{2}$ 以外的范畴 --- 比如那些含有非平凡态射的范畴 --- 就会对 $\Delta_c$
和 $D$ 之间的转换施加自然性条件. 我们把这样的转换叫作\emph{锥}, 因为 $\Delta$ 的像是锥/棱锥的
顶点, 锥的侧面由自然转换的分量构成. $D$ 的像构成锥的底.

一般说来, 要构造一个锥, 我们从一个定义模式的范畴 $\cat{I}$ 开始. 它是个小范畴, 常常是有限范畴.
我们从 $\cat{I}$ 到 $\cat{C}$ 挑一个函子 $D$, 把它 (或它的像) 叫作\emph{图式}. 我们在
$\cat{C}$ 里挑某个 $c$ 作为锥的顶点, 用它来定义从 $\cat{I}$ 到 $\cat{C}$ 的常数函子
$\Delta_c$. 从 $\Delta_c$ 到 $D$ 的一个自然转换就是我们的锥. 对有限的 $\cat{I}$ 而言, 它不过是
一堆把 $c$ 连到图式 --- $\cat{I}$ 在 $D$ 下的像 --- 的态射.

\begin{figure}[H]
  \centering
  \includegraphics[width=0.35\textwidth]{images/cone.jpg}
\end{figure}

\noindent
自然性要求这个图式里的所有三角形 (棱锥的侧面) 都交换. 事实上, 取 $\cat{I}$ 中任一态射 $f$.
函子 $D$ 把它映成 $\cat{C}$ 中的一个态射 $D f$, 它构成某个三角形的底边. 常数函子 $\Delta_c$
把 $f$ 映成 $c$ 上的恒等态射. $\Delta$ 把态射的两端挤成一个物件, 于是自然性方形就成了一个交换
三角形. 这个三角形的两条腰就是自然转换的分量.

\begin{figure}[H]
  \centering
  \includegraphics[width=0.35\textwidth]{images/conenaturality.jpg}
\end{figure}

\noindent
这就是一个锥. 我们关心的是\newterm{universal cone 普遍锥} --- 就像我们定义积时挑了一个普遍物件一样.

做到这一点有很多办法. 比如, 我们可以基于给定的函子 $D$ 定义一个\emph{锥的范畴}. 那个范畴里的
物件就是锥. 不过, $\cat{C}$ 中并非每个物件 $c$ 都能当锥的顶点, 因为 $\Delta_c$ 和 $D$ 之间
可能根本没有自然转换.

要让它成为一个范畴, 还得定义锥之间的态射. 这些态射应当完全由它们顶点之间的态射决定. 但不是随便哪个
态射都行. 记住, 在构造积的时候我们施加过一个条件: 候选物件 (也就是顶点) 之间的态射必须是各投影的
公因子. 例如:

\src{snippet01}

\begin{figure}[H]
  \centering
  \includegraphics[width=0.35\textwidth]{images/productranking.jpg}
\end{figure}

推广开来, 这个条件就变成: 以分解态射为一条边的那些三角形全都交换.

\begin{figure}[H]
  \centering
  \includegraphics[width=0.35\textwidth]{images/conecommutativity.jpg}
  \caption{连接两个锥的交换三角形, 其中 $h$ 是分解态射 (这里, 下方那个锥是普遍锥, 顶点为
    $\Lim[D]$)}
\end{figure}

\noindent
我们就把这些分解态射当作锥的范畴里的态射. 很容易验证它们确实可以组合, 而且恒等态射也是一个分解态射.
于是锥构成一个范畴.

现在我们可以把普遍锥定义为锥的范畴里的\emph{终止物件}. 终止物件的定义说, 从任何其他物件到它都有
唯一的态射. 在我们这里, 这意味着从任意其他锥的顶点到普遍锥的顶点有唯一的分解态射. 我们把这一普遍锥
叫作图式 $D$ 的\emph{极限}, $\Lim[D]$ (文献里你常会看到 $\Lim$ 符号下面有个指向 $I$ 的左箭头).
通常我们会偷懒, 直接把这个锥的顶点叫作极限 (或极限物件).

直觉上, 极限把整个图式的性质体现在了单个物件之中. 比如, 我们那个两物件图式的极限就是两个物件的积.
积 (连同两个投影) 包含了关于这两个物件的信息. 而普遍性则意味着它不含多余的垃圾.

\section{极限作为自然同构}

这个极限的定义仍有令人不满意的地方. 我是说, 它管用, 但我们还带着那个交换性条件 --- 用来联系任意
两个锥的那些三角形. 如果能把它换成某个自然性条件, 那就优雅多了. 可怎么做?

我们要处理的不再是一个锥, 而是整整一簇锥 (实际上是一个范畴). 如果极限存在 (把话说清楚 --- 并不保证
它一定存在), 那么这些锥中就有一个是普遍锥. 对每个其他的锥, 我们都有一个唯一的分解态射, 把它的顶点
 (就叫它 $c$ 吧) 映到普遍锥的顶点, 我们把这个顶点命名为 $\Lim[D]$. (其实我可以省掉 ``其他''
这个词, 因为恒等态射把普遍锥映到自己, 而它平凡地经由自身分解.) 让我把要紧的部分重复一遍: 给定任意
一个锥, 都有唯一一个特殊种类的态射. 我们有一个从锥到特殊态射的映射, 而且它是一一映射.

这个特殊的态射是 hom-集 $\cat{C}(c, \Lim[D])$ 的成员. 这个 hom-集里的其他成员就没那么幸运了,
因为它们无法分解两个锥之间的映射. 我们想要的是: 对每个 $c$, 都能从集合 $\cat{C}(c, \Lim[D])$
中挑出一个态射 --- 一个满足那个特定交换性条件的态射. 这听起来像是在定义自然转换吗? 确确实实是!

可这个转换联系的是哪两个函子?

一个函子是把 $c$ 映到集合 $\cat{C}(c, \Lim[D])$ 的这个映射. 它是从 $\cat{C}$ 到 $\Set$ 的
函子 --- 它把物件映到集合. 实际上它是个逆变函子. 我们这样定义它对态射的作用: 取一个从 $c'$ 到
$c$ 的态射 $f$:
\[f \Colon c' \to c\]
我们的函子把 $c'$ 映到集合 $\cat{C}(c', \Lim[D])$. 要定义这个函子对 $f$ 的作用 (换句话说,
要提升 $f$), 我们必须定义 $\cat{C}(c, \Lim[D])$ 和 $\cat{C}(c', \Lim[D])$ 之间的相应映射.
那我们就取 $\cat{C}(c, \Lim[D])$ 的一个元素 $u$, 看看能不能把它映到
$\cat{C}(c', \Lim[D])$ 的某个元素. hom-集的元素是一个态射, 所以我们有:
\[u \Colon c \to \Lim[D]\]
我们可以把 $u$ 与 $f$ 预组合, 得到:
\[u . f \Colon c' \to \Lim[D]\]
而这是 $\cat{C}(c', \Lim[D])$ 的一个元素 --- 所以没错, 我们找到了一个态射之间的映射:

\src{snippet02}
注意 $c$ 和 $c'$ 顺序的颠倒, 这是\emph{逆变}函子的特征.

\begin{figure}[H]
  \centering
  \includegraphics[width=0.35\textwidth]{images/homsetmapping.jpg}
\end{figure}

\noindent
要定义一个自然转换, 我们还需要另一个函子, 它也是从 $\cat{C}$ 到 $\Set$ 的映射. 但这一次我们要
考虑一组锥. 锥不过是自然转换, 所以我们看的是自然转换的集合 $Nat(\Delta_c, D)$. 从 \code{c}
到这个特定的自然转换集合的映射是一个 (逆变) 函子. 我们怎么证明这点? 同样, 来定义它对态射的作用:
\[f \Colon c' \to c\]
$f$ 的提升应当是一个自然转换之间的映射, 而这些自然转换介于两个从 $\cat{I}$ 到 $\cat{C}$ 的函子
之间:
\[Nat(\Delta_c, D) \to Nat(\Delta_{c'}, D)\]
我们怎么映射自然转换? 每个自然转换都是一次态射的挑选 --- 它的分量 --- $\cat{I}$ 的每个元素对应
一个态射. 某个 $\alpha$ (它是 $Nat(\Delta_c, D)$ 的成员) 在 $a$ (它是 $\cat{I}$ 中的一个物件)
处的分量是一个态射:
\[\alpha_a \Colon \Delta_c a \to D a\]
或者, 用常数函子 $\Delta$ 的定义写成,
\[\alpha_a \Colon c \to D a\]
给定 $f$ 和 $\alpha$, 我们必须构造一个 $\beta$, 它是 $Nat(\Delta_{c'}, D)$ 的成员. 它在
$a$ 处的分量应当是一个态射:
\[\beta_a \Colon c' \to D a\]
我们可以很容易地由前者 ($\alpha_a$) 得到后者 ($\beta_a$), 只要把它与 $f$ 预组合:
\[\beta_a = \alpha_a . f\]
证明这些分量确实汇成了一个自然转换相对容易.

\begin{figure}[H]
  \centering
  \includegraphics[width=0.4\textwidth]{images/natmapping.jpg}
\end{figure}

\noindent
给定态射 $f$, 我们就此逐分量地建立了两个自然转换之间的映射. 这个映射为该函子定义了
\code{contramap}:
\[c \to Nat(\Delta_c, D)\]
我刚才做的事情, 是向你表明我们有两个从 $\cat{C}$ 到 $\Set$ 的 (逆变) 函子. 我没有做任何假设
--- 这两个函子总是存在.

顺带一提, 这两个函子中的第一个在范畴论里扮演着重要角色, 等我们讲米田引理时还会再见到它. 从任意
范畴 $\cat{C}$ 到 $\Set$ 的逆变函子有个名字: 它们叫作 ``预层''. 这里这个叫\newterm{representable
presheaf 可表预层}. 第二个函子也是个预层.

既然有了两个函子, 我们就能谈它们之间的自然转换. 那就不多客套, 直接给出结论: 从 $\cat{I}$ 到
$\cat{C}$ 的函子 $D$ 拥有极限 $\Lim[D]$, 当且仅当我刚定义的那两个函子之间存在一个自然同构:
\[\cat{C}(c, \Lim[D]) \simeq Nat(\Delta_c, D)\]
提醒一下什么是自然同构. 它是一个每个分量都是同构的自然转换, 也就是说, 每个分量都是可逆态射.

我不打算把这个命题的证明走一遍. 证明过程即便不繁琐, 也相当直截了当. 处理自然转换时, 你通常把注意力
放在分量上, 它们是态射. 在本例中, 由于两个函子的目标都是 $\Set$, 自然同构的分量将是函数. 这些是
高阶函数, 因为它们从 hom-集映到自然转换的集合. 同样, 你可以通过考察一个函数对参数做了什么来分析它:
这里的参数是一个态射 --- $\cat{C}(c, \Lim[D])$ 的成员 --- 结果是一个自然转换 ---
$Nat(\Delta_c, D)$ 的成员, 也就是我们所说的锥. 而这个自然转换本身又有它自己的分量, 它们也是
态射. 所以一层层下去全是态射, 只要能盯住它们, 你就能证明这个命题.

最重要的结果是: 这个同构的自然性条件, 恰好就是锥的映射的那个交换性条件.

预告一下后面的内容: 集合 $Nat(\Delta_c, D)$ 可以看成函子范畴中的一个 hom-集; 于是我们的自然同构
联系了两个 hom-集, 这就指向了一个更加普遍的关系, 叫作伴随.

\section{极限的例子}

我们已经看到, 范畴论中的积是由一个我们叫作 $\cat{2}$ 的简单范畴所生成的图式的极限.

还有一个更简单的极限例子: 终止物件. 第一反应会以为是单例范畴导出了终止物件, 但真相比那更极端:
终止物件是由空范畴生成的极限. 从空范畴出发的函子不挑选任何物件, 于是锥收缩得只剩下顶点. 普遍锥就
是那个孤零零的顶点, 从任何其他顶点都有唯一的态射指向它. 你会发现, 这正是终止物件的定义.

下一个有趣的极限叫\emph{等化子}. 它是由一个两物件范畴生成的极限, 两个物件之间有两个平行态射
 (以及, 照例, 恒等态射). 这个范畴挑选出的 $\cat{C}$ 中的图式由两个物件 $a$ 和 $b$
以及两个态射构成:

\src{snippet03}

要在这一图式上构造锥, 我们必须加上顶点 $c$ 和两个投影:

\src{snippet04}

\begin{figure}[H]
  \centering
  \includegraphics[width=0.35\textwidth]{images/equalizercone.jpg}
\end{figure}

\noindent
我们有两个三角形必须交换:

\src{snippet05}

这告诉我们, $q$ 由这些方程中的某一个唯一确定, 比如说 \code{q = f . p}, 于是我们可以把它从图里
省掉. 所以我们只剩下一个条件:

\src{snippet06}

理解它的方式是: 如果我们把注意力限制在 $\Set$ 上, 函数 $p$ 的像挑选出 $a$ 的一个子集. 限制
在这个子集上, 函数 $f$ 和 $g$ 相等.

比如, 取 $a$ 为由坐标 $x$ 和 $y$ 参数化的二维平面. 取 $b$ 为实数直线, 并取:

\src{snippet07}

这两个函数的等化子是实数集 (也就是顶点 $c$) 加上函数:

\src{snippet08}

注意 $(p~t)$ 定义了二维平面中的一条直线. 沿着这条直线, 两个函数相等.

当然, 还有别的集合 $c'$ 和函数 $p'$ 也能导向这个等式:

\src{snippet09}

但它们全都唯一地经由 $p$ 分解出来. 例如, 我们可以取单例集合 $\cat{()}$ 作为 $c'$ 和函数:

\src{snippet10}

它是个不错的锥, 因为 $f (0, 0) = g (0, 0)$. 但它不普遍, 因为存在经由 $h$ 的唯一分解:

\src{snippet11}

其中

\src{snippet12}

\begin{figure}[H]
  \centering
  \includegraphics[width=0.35\textwidth]{images/equilizerlimit.jpg}
\end{figure}

\noindent
等化子因此可以用来求解 $f~x = g~x$ 类型的方程. 但它远比这更一般, 因为它是用物件和态射而不是用
代数方式定义的.

求解方程这一思想更进一步, 体现在另一个极限之中 --- 拉回. 这里我们依然想等同两个态射, 但这一次它们的
定义域不同. 我们从一个形状为 $1\rightarrow2\leftarrow3$ 的三物件范畴开始. 对应于这个范畴
的图式由三个物件 $a$, $b$, $c$ 和两个态射构成:

\src{snippet13}

这个图式常被称为\emph{cospan}.

在这个图式上搭建的锥由顶点 $d$ 和三个态射构成:

\src{snippet14}

\begin{figure}[H]
  \centering
  \includegraphics[width=0.35\textwidth]{images/pullbackcone.jpg}
\end{figure}

\noindent
交换性条件告诉我们, $r$ 完全由其余态射决定, 可以从图里省掉. 所以我们只剩下下面这个条件:

\src{snippet15}
拉回就是这种形状的普遍锥.

\begin{figure}[H]
  \centering
  \includegraphics[width=0.35\textwidth]{images/pullbacklimit.jpg}
\end{figure}

\noindent
同样, 如果你把视野收窄到集合, 就可以把物件 $d$ 想成由 $a$ 和 $c$ 中的元素组成的对, 其中
$f$ 作用于第一个分量的结果等于 $g$ 作用于第二个分量的结果. 如果这还太一般, 不妨考虑 $g$ 是
常数函数的特例, 比如 $g~\_ = 1.23$ (假定 $b$ 是实数集合). 那么你真的就是在求解方程:

\src{snippet16}

这种情况下, $c$ 的选择无关紧要 (只要它不是空集), 所以我们可以把它取成单例集合.
$a$ 可以是三维向量的集合, 而 $f$ 是向量长度. 那么拉回就是所有对 $(v, ())$ 的集合,
其中 $v$ 是长度为 1.23 的向量 (方程 $\sqrt{(x^{2}+y^{2}+z^{2})} = 1.23$ 的解), 而 $()$ 是
单例集合中的哑元素.

但拉回还有更一般的应用, 在编程中也是. 比如, 把 C++ 的类看成一个范畴, 其中的态射是把子类连到父类的
箭头. 我们把继承当作传递性质, 所以如果 \code{C} 继承自 \code{B} 而 \code{B} 继承自 \code{A},
我们就说 \code{C} 继承自 \code{A} (毕竟, 你可以把指向 \code{C} 的指针当作指向 \code{A} 的指针来
传). 而且我们假定 \code{C} 继承自 \code{C}, 于是每个类都有恒等箭头. 这样一来, 子类化就和子类型
对齐了. C++ 还支持多重继承, 于是你可以构造一个菱形继承图式: 两个类 \code{B} 和 \code{C} 继承自
\code{A}, 第四个类 \code{D} 多重继承自 \code{B} 和 \code{C}. 通常 \code{D} 会得到两份 \code{A}
的拷贝, 这很少是想要的; 但你可以用虚继承让 \code{D} 中只有一份 \code{A}.

让 \code{D} 成为这个图式中的拉回意味着什么? 那意味着任何多重继承自 \code{B} 和 \code{C} 的类
\code{E} 也都是 \code{D} 的子类. 这在 C++ 里无法直接表达, 因为 C++ 的子类型是名义上的 (C++ 编译器
不会推断出这种类关系 --- 那需要 ``鸭子类型''). 但我们可以走到子类型关系之外, 转而询问从 \code{E} 到
\code{D} 的转换是否安全. 如果 \code{D} 只是 \code{B} 和 \code{C} 的光板组合, 没有额外数据,
也没有覆盖任何方法, 那么这个转换就是安全的. 当然, 如果 \code{B} 和 \code{C} 的某些方法之间存在名字
冲突, 拉回就不存在.

\begin{figure}[H]
  \centering
  \includegraphics[width=0.25\textwidth]{images/classes.jpg}
\end{figure}

\noindent
拉回在类型推断中还有一个更高级的用法. 我们常常需要\emph{统一}两个表达式的类型. 比如, 假设编译器要推断
一个函数的类型:

\begin{snip}{haskell}
twice f x = f (f x)
\end{snip}
它会给所有变量和子表达式赋上初步类型. 特别地, 它会赋:

\begin{snip}{haskell}
f       :: t0
x       :: t1
f x     :: t2
f (f x) :: t3
\end{snip}
从中它将推出:

\begin{snip}{haskell}
twice :: t0 -> t1 -> t3
\end{snip}
它还会得到一组由函数应用规则导出的约束:

\begin{snip}{haskell}
t0 = t1 -> t2 -- because f is applied to x 
t0 = t2 -> t3 -- because f is applied to (f x)
\end{snip}
这些约束必须靠找到一组类型 (或类型变量) 来统一, 它们代入两个表达式中的未知类型后能产生相同的类型.
一个这样的代入是:

\begin{snip}{haskell}
t1 = t2 = t3 = Int 
twice :: (Int -> Int) -> Int -> Int
\end{snip}
但显然它不是最一般的代入. 最一般的代入要靠拉回才能得到. 我不打算展开细节, 因为那超出了本书范围, 但你
可以自己确认结果应该是:

\begin{snip}{haskell}
twice :: (t -> t) -> t -> t
\end{snip}
其中 \code{t} 是一个自由类型变量.

\section{余极限}

就像范畴论中的一切构造一样, 极限在反范畴中有它的对偶影像. 当你把锥里所有箭头的方向颠倒过来,
就得到一个余锥, 而那些余锥中的普遍者叫作余极限. 注意这个颠倒也影响了分解态射, 它现在从普遍余锥
流向任意其他余锥.

\begin{figure}[H]
  \centering
  \includegraphics[width=0.35\textwidth]{images/colimit.jpg}
  \caption{带分解态射 $h$ 的余锥, 连接两个顶点.}
\end{figure}

\noindent
余极限的一个典型例子是余积, 它对应由于 $\cat{2}$ 生成的图式 --- 我们在积的定义里用过这个范畴.

\begin{figure}[H]
  \centering
  \includegraphics[width=0.35\textwidth]{images/coproductranking.jpg}
\end{figure}

\noindent
积和余积都以各自不同的方式体现了一对物件的本质.

正如终止物件是个极限, 初始物件也是个余极限, 对应于基于空范畴的图式.

拉回的对偶叫\emph{推出}. 它基于一个叫 span 的图式, 由范畴 $1\leftarrow2\rightarrow3$
生成.

\section{连续性}

我之前说过, 函子接近于范畴的连续映射这一想法, 因为它们从不拆掉已有的连接 (态射). 从一个范畴
$\cat{C}$ 到 $\cat{C'}$ 的\emph{连续函子} $F$ 的实际定义, 还额外要求函子保持极限. 每个在
$\cat{C}$ 中的图式 $D$ 都可以映成 $\cat{C'}$ 中的图式 $F \circ D$, 只要把两个函子组合起来.
$F$ 的连续性条件说: 如果图式 $D$ 有极限 $\Lim[D]$, 那么图式 $F \circ D$ 也有极限, 而且它等于
$F (\Lim[D])$.

\begin{figure}[H]
  \centering
  \includegraphics[width=0.6\textwidth]{images/continuity.jpg}
\end{figure}

\noindent
注意, 因为函子把态射映到态射, 把组合映到组合, 所以锥的像总是个锥. 交换三角形总是被映成交换三角形
(函子保持组合). 分解态射也一样: 分解态射的也是个分解态射. 所以每个函子都\emph{差点儿}就连续了.
可能出问题的是唯一性条件. $\cat{C'}$ 中的分解态射可能不是唯一的. $\cat{C'}$ 中也可能存在其他
``更好的锥'', 它们在 $\cat{C}$ 中并不存在.

hom-函子是连续函子的一个例子. 回顾一下, hom-函子 $\cat{C}(a, b)$ 对第一个变量逆变, 对第二个
变量协变. 换句话说, 它是一个函子:
\[\cat{C}^{op} \times \cat{C} \to \Set\]
当它的第二个参数固定时, hom-集函子 (它就成了可表预层) 把 $\cat{C}$ 中的余极限映到 $\Set$ 中的
极限; 而当它的第一个参数固定时, 它把极限映到极限.

在 Haskell 里, hom-函子就是把任意两个类型映到一个函数类型, 所以它只是一个参数化的函数类型. 当我们
固定第二个参数, 比如说固定为 \code{String}, 得列逆变函子:

\src{snippet17}
连续性意味着, 当 \code{ToString} 作用于一个余极限时, 比如余积 \code{Either b c}, 它会产出
一个极限; 在本例中是两个函数类型的积:

\src{snippet18}
确实, 任何以 \code{Either b c} 为参数的函数都实现为一个 case 语句, 两个 case 分给一对函数
处理.

类似地, 当我们固定 hom-集的第一个参数, 就得到熟悉的读者函子. 它的连续性意味着, 比如, 任何返回积的
函数等价于一对函数; 特别地:

\src{snippet19}
我知道你在想什么: 用不着范畴论你也能看出这些东西. 你说得对! 不过我仍然觉得不可思议: 这样的结果可以从
第一性原理推导出来, 而无需涉及位和字节、处理器架构、编译器技术, 甚至无需 lambda 演算.

如果你好奇 ``极限'' 和 ``连续性'' 这些名字从何而来, 它们是微积分中相应概念的推广. 在微积分里, 极限和
连续性是用开邻域定义的. 开集 (用来定义拓扑) 构成一个范畴 (一个偏序集).

\section{挑战}

\begin{enumerate}
  \tightlist
  \item
        你会如何描述 C++ 类范畴中的推出?
  \item
        证明恒等函子 $\mathbf{Id} \Colon \cat{C} \to \cat{C}$ 的极限是初始物件.
  \item
        给定集合的子集构成一个范畴. 该范畴中的态射定义为: 当第一个集合是第二个集合的子集时,
        连接两个集合的箭头. 在这样的范畴里, 两个集合的拉回是什么? 推出是什么?
        初始物件和终止物件又是什么?
  \item
        你能猜到余等化子是什么吗?
  \item
        证明, 在一个有终止物件的范畴中, 朝终止物件的拉回就是积.
  \item
        类似地, 证明从初始物件出发的推出 (如果它存在) 就是余积.
\end{enumerate}
